Now it follows from 2.2 that the isomorphism:
@@M0@@
defined by the isomorphisms (2.5.2) is compatible with the algebra structures of the two sides; hence the same is true of the isomorphism:
@@M1@@
defined by the isomorphisms (2.5.3). By uniqueness, the latter is therefore the restriction of the homomorphism (2.5.1) that was defined globally. The homomorphism (2.5.1) is thus an isomorphism.

ii) Consider the exact sequence:
@@M2@@
and the resulting exact @@M3@@ sequence:
@@M4@@
It follows that @@M5@@ is the kernel of the homomorphism:
@@M6@@
that is, @@M7@@.

To prove the last relation, we shall show that there is an exact sequence:
@@M8@@
The homomorphisms @@M9@@ and @@M10@@ were defined in (1.6.1); here we take their difference. Furthermore, @@M11@@ maps naturally into @@M12@@, hence into @@M13@@, and likewise into @@M14@@, and these homomorphisms vanish on @@M15@@, since @@M16@@ (where @@M17@@ denotes the ideal of @@M18@@), defining the arrow on the left.
